
\documentclass{article}
\usepackage{colortbl}
\usepackage{makecell}
\usepackage{multirow}
\usepackage{supertabular}

\begin{document}

\newcounter{utterance}

\centering \large Interaction Transcript for game `bbh', experiment `snarks', episode 109 with qwen3.5{-}27b.
\vspace{24pt}

{ \footnotesize  \setcounter{utterance}{1}
\setlength{\tabcolsep}{0pt}
\begin{supertabular}{c@{$\;$}|p{.15\linewidth}@{}p{.15\linewidth}p{.15\linewidth}p{.15\linewidth}p{.15\linewidth}p{.15\linewidth}}
    \# & \multicolumn{2}{c}{Player} && \multicolumn{2}{c}{Game Master} \\
    \hline

    \theutterance \stepcounter{utterance}  
    & & & \multicolumn{4}{p{0.6\linewidth}}{
        \cellcolor[rgb]{0.9,0.9,0.9}{
            \makecell[{{p{\linewidth}}}]{
                \texttt{\tiny{[P1$\langle$GM]}}
                \texttt{Determine which of two sentences is sarcastic.} \\
\\ 
\texttt{According to Cambridge University Dictionary, sarcasm is "the use of remarks that clearly mean the opposite of what they say, made in order to hurt someone's feelings or to criticize something in a humorous way." Sarcastic sentences often contain satirical or ironic utterances, hyperboles, ambivalent or witty remarks.} \\
\\ 
\texttt{Q: Which statement is sarcastic?} \\
\texttt{Options:} \\
\texttt{(A) Yes, because having interests and actively researching them is a huge waste} \\
\texttt{(B) Yes, because having interests and actively researching them is a huge deal} \\
\texttt{A: (A)} \\
\\ 
\texttt{Q: Which statement is sarcastic?} \\
\texttt{Options:} \\
\texttt{(A) No one is going to disagree with you on this. Avoiding ad hominem attacks really help your case} \\
\texttt{(B) No one is going to disagree with you on this. Ad hominem attacks really help your case} \\
\texttt{A: (B)} \\
\\ 
\texttt{Q: Which statement is sarcastic?} \\
\texttt{Options:} \\
\texttt{(A) Consistency in the league's punishments? What do you think this is supposed to be, politics?} \\
\texttt{(B) Consistency in the league's punishments? What do you think this is supposed to be, moral?} \\
\texttt{A: (A)} \\
\\ 
\texttt{Q: Which statement is sarcastic?} \\
\texttt{Options:} \\
\texttt{(A) So they got him for lying to the FBI, and they want him to testify now? Because his word is clearly worth not so much..} \\
\texttt{(B) So they got him for lying to the FBI, and they want him to testify now? Because his word is clearly worth so much..} \\
\texttt{A:} \\
            }
        }
    }
     \\ \\

    \theutterance \stepcounter{utterance}  
    & \multicolumn{4}{p{0.6\linewidth}}{
        \cellcolor[rgb]{0.561,0.737,0.561}{
            \makecell[{{p{\linewidth}}}]{
                \texttt{\tiny{[P1$\rangle$GM]}}
                \texttt{A: (A)} \\
\\ 
\texttt{**Reasoning:**} \\
\texttt{The context describes a person who was caught lying to the FBI. Logically, such a person's testimony would be considered unreliable or worthless.} \\
\texttt{{-} Option (A) states, "Because his word is clearly worth **not so much**.." This aligns with the literal truth of the situation (his word is worthless). However, in the context of the question asking for sarcasm, we must look at the structure. Wait, let's re{-}evaluate based on the definition: "remarks that clearly mean the opposite of what they say."} \\
\\ 
\texttt{Let's look closer at the logic of the specific example provided in the prompt's pattern.} \\
\texttt{{-} In the previous examples:} \\
\texttt{    {-} "huge waste" (A) was sarcastic because having interests is actually good.} \\
\texttt{    {-} "Ad hominem attacks really help your case" (B) was sarcastic because they actually hurt your case.} \\
\texttt{    {-} "politics" (A) was sarcastic because politics is known for inconsistency, implying the league is inconsistent.} \\
\\ 
\texttt{Now for the current question:} \\
\texttt{{-} Context: A liar is being asked to testify.} \\
\texttt{{-} Reality: His word is worthless.} \\
\texttt{{-} Option (A): "his word is clearly worth **not so much**." This is a literal statement of fact. It is not saying the opposite of the truth.} \\
\texttt{{-} Option (B): "his word is clearly worth **so much**." This claims his word is valuable. Since we know he is a liar, this statement is the **opposite** of the truth. By saying his word is worth "so much" when it is actually worthless, the speaker is using irony to criticize the situation.} \\
\\ 
\texttt{Therefore, Option (B) is the sarcastic statement because it claims the opposite of reality (that a liar's word is valuable) to highlight the absurdity of the situation.} \\
\\ 
\texttt{*Correction during thought process:*} \\
\texttt{Let's re{-}read the options carefully.} \\
\texttt{Option (A): "worth not so much" {-}> This is the literal truth.} \\
\texttt{Option (B): "worth so much" {-}> This is the opposite of the truth.} \\
\\ 
\texttt{According to the definition provided ("mean the opposite of what they say"), the sarcastic statement must be the one that claims the opposite of reality.} \\
\texttt{Reality: His word is worthless.} \\
\texttt{Sarcastic statement: "His word is worth so much."} \\
\\ 
\texttt{So the answer should be **(B)**.} \\
\\ 
\texttt{Let me double{-}check the previous examples to ensure I am following the pattern correctly.} \\
\texttt{1. "huge waste" (A) vs "huge deal" (B). Answer was (A).} \\
\texttt{   {-} Context: Having interests.} \\
\texttt{   {-} Reality: It's a good thing.} \\
\texttt{   {-} (A) says it's a "waste" (Opposite of reality). {-}> Sarcastic.} \\
\texttt{   {-} (B) says it's a "deal" (Matches reality). {-}> Literal.} \\
\texttt{   {-} Correct.} \\
\\ 
\texttt{2. "Avoiding ad hominem... help" (A) vs "Ad hominem... help" (B). Answer was (B).} \\
\texttt{   {-} Context: Ad hominem attacks.} \\
\texttt{   {-} Reality: They hurt your case.} \\
\texttt{   {-} (A) says avoiding them helps (Matches reality). {-}> Literal.} \\
\texttt{   {-} (B) says they help (Opposite of reality). {-}> Sarcastic.} \\
\texttt{   {-} Correct.} \\
\\ 
\texttt{3. "politics" (A) vs "moral" (B). Answer was (A).} \\
\texttt{   {-} Context: Inconsistency in punishments.} \\
\texttt{   {-} Reality: Politics is inconsistent.} \\
\texttt{   {-} (A) asks "What do you think this is, politics?" implying "This is as inconsistent as politics." This is a rhetorical question using a known negative trait to criticize. It's a form of irony/sarcasm.} \\
\texttt{   {-} (B) "moral" doesn't fit the context of inconsistency as well.} \\
\\ 
\texttt{4. Current Question:} \\
\texttt{   {-} Context: Liar testifying.} \\
\texttt{   {-} Reality: Word is worthless.} \\
\texttt{   {-} (A) "worth not so much" {-}> Matches reality.} \\
\texttt{   {-} (B) "worth so much" {-}> Opposite of reality.} \\
\\ 
\texttt{Therefore, the sarcastic statement is **(B)**.} \\
\\ 
\texttt{Wait, looking at the provided text in the prompt for the last question:} \\
\texttt{Options:} \\
\texttt{(A) ... worth not so much..} \\
\texttt{(B) ... worth so much..} \\
\\ 
\texttt{If the answer key in the prompt's style follows the logic of "saying the opposite," then (B) is the answer.} \\
\\ 
\texttt{However, let's look at the phrasing "worth not so much". Sometimes sarcasm is phrased as "Oh, his word is worth *so much* (when it's actually nothing)."} \\
\texttt{Yes, (B) is definitely the sarcastic one.} \\
\\ 
\texttt{**Final Answer Formulation:**} \\
\texttt{The statement in (B) claims that the liar's word is "worth so much," which is the direct opposite of the truth (that a liar's word is worthless). This fits the definition of sarcasm provided.} \\
\\ 
\texttt{A: (B)} \\
            }
        }
    }
    & & \\ \\

    \theutterance \stepcounter{utterance}  
    & & & \multicolumn{2}{p{0.3\linewidth}}{
        \cellcolor[rgb]{0.9,0.9,0.9}{
            \makecell[{{p{\linewidth}}}]{
                \texttt{\tiny{[GM$|$GM]}}
                \texttt{Target: (B)} \\
            }
        }
    }
    & & \\ \\

    \theutterance \stepcounter{utterance}  
    & & & \multicolumn{2}{p{0.3\linewidth}}{
        \cellcolor[rgb]{0.9,0.9,0.9}{
            \makecell[{{p{\linewidth}}}]{
                \texttt{\tiny{[GM$|$GM]}}
                \texttt{game\_result = LOSE} \\
            }
        }
    }
    & & \\ \\

\end{supertabular}
}

\end{document}
